\documentclass[12pt,a4paper]{article}

% Required packages
\usepackage[left=4cm,right=3cm,top=3cm,bottom=3cm]{geometry} % Margins
\usepackage{times}               % Times New Roman font
\usepackage[utf8]{inputenc}      % Input encoding for UTF-8 characters
\usepackage[T1]{fontenc}         % Font encoding for proper character rendering
\usepackage[indonesian]{babel}   % Indonesian language support
\usepackage{setspace}            % For line spacing
\usepackage{ragged2e}            % For text justification
\usepackage{tocloft}             % For table of contents formatting
\usepackage{titlesec}            % For section formatting
\usepackage{lipsum}              % For lorem ipsum
\usepackage[hidelinks]{hyperref} % For clickable links in TOC
\usepackage{graphicx}            % Package for handling graphics/images
\usepackage{scrlayer-scrpage}    % Package for header and footer customization
\usepackage[hidelinks]{hyperref} % Package for clickable links without visible borders
\usepackage{tabularx}            % Package for advanced table formatting with auto-adjusting columns
\usepackage{makecell}            % Package for creating cells with line breaks in tables
\usepackage[none]{hyphenat}      % Package to disable hyphenation
\usepackage{xstring}             % Package for string manipulation and operations
\usepackage{pgffor}              % Package for foreach loops and iterations
\usepackage{enumitem}            % Package for customizing lists and enumerations
% \usepackage{showframe}

% Document settings
\setstretch{1.15}            % 1.15 line spacing
% \justifying                  % Justified text alignment
\newcommand{\judul}{<Judul Kita di Sini>}

% Dosen Pendamping Variables
\newcommand{\dosenNama}{Prof. Dr. Andre Indrawan, M.Hum., M.Mus.}
\newcommand{\dosenGender}{Laki-laki}
\newcommand{\dosenProdi}{Musik (S1)}
\newcommand{\dosenNIP}{196105101987031002}
\newcommand{\dosenNIDN}{0010056110}
\newcommand{\dosenTTL}{Bandung, 10 Mei 1961}
\newcommand{\dosenEmail}{indrawan\_andre@isi.ac.id}
\newcommand{\dosenTelp}{+628180-4251709}

% Ketua Kelompok Variables
\newcommand{\ketuaNama}{Vincent Nuridzati Adittama}
\newcommand{\ketuaGender}{Laki-laki}
\newcommand{\ketuaProdi}{Musik (S1)}
\newcommand{\ketuaNIM}{23104810131}
\newcommand{\ketuaTTL}{Trenggalek, 23 November 2004}
\newcommand{\ketuaEmail}{eezaateedanvers@gmail.com}
\newcommand{\ketuaTelp}{+62 822 4697 9795}

% Anggota 1 Variables
\newcommand{\anggotaSatuNama}{Azzahra Qurrata A'yun Adindya Irbah Ramadhani}
\newcommand{\anggotaSatuGender}{Perempuan}
\newcommand{\anggotaSatuProdi}{Musik (S1)}
\newcommand{\anggotaSatuNIM}{24106110131}
\newcommand{\anggotaSatuTTL}{Wonosobo, 7 Oktober 2004}
\newcommand{\anggotaSatuEmail}{irbahramadhani7@gmail.com}
\newcommand{\anggotaSatuTelp}{+62 857 4793 0265}

% Anggota 2 Variables
\newcommand{\anggotaDuaNama}{Elizabeth Ardhayu Maheswari}
\newcommand{\anggotaDuaGender}{Perempuan}
\newcommand{\anggotaDuaProdi}{Musik}
\newcommand{\anggotaDuaNIM}{23104330131}
\newcommand{\anggotaDuaTTL}{Lampung, 26 September 2005}
\newcommand{\anggotaDuaEmail}{ardhayu.maheswari@gmail.com}
\newcommand{\anggotaDuaTelp}{+62 813 2964 8109}

% Anggota 3 Variables
\newcommand{\anggotaTigaNama}{Cherish Coo Shalom Putri Betzia}
\newcommand{\anggotaTigaGender}{Perempuan}
\newcommand{\anggotaTigaProdi}{Musik (S1)}
\newcommand{\anggotaTigaNIM}{24106080131}
\newcommand{\anggotaTigaTTL}{Yogyakarta, 17 Agustus 2006}
\newcommand{\anggotaTigaEmail}{cherishalom07@gmail.com}
\newcommand{\anggotaTigaTelp}{+62 878 1808 7178}

% Section formatting
\addto\captionsindonesian{
  \renewcommand{\contentsname}{DAFTAR ISI}
  \renewcommand{\refname}{DAFTAR PUSTAKA}
  \renewcommand{\abstractname}{ABSTRAK}
  \renewcommand{\figurename}{Gambar}
  \renewcommand{\tablename}{Tabel}
}

% Section formatting
\titleformat{\section}{\large\normalfont\bfseries}{BAB \thesection.}{1em}{}
\titleformat{\subsection}{\normalsize\normalfont\bfseries}{\thesubsection.}{1em}{}

% Page numbering and header/footer settings
\clearpairofpagestyles
\cfoot[]{}

% For front matter (Table of Contents)
\newpairofpagestyles{frontmatter}{
    \KOMAoptions{footsepline=false}
    \cfoot[]{}
    \ofoot*{\pagemark}
}

% For main matter (Chapter 1 onwards)
\newpairofpagestyles{mainmatter}{
    \KOMAoptions{headsepline=false}
    \chead[]{}
    \ohead*{\pagemark}
}

% Name Shortening to Initials after 3 words
\newcount\wordcount
\newcommand{\formatName}[1]{%
  \wordcount=0
  \def\result{}%
  \expandafter\parsewords#1 \relax\relax
  \result
}

\def\parsewords#1 #2\relax{%
  \ifx\relax#1\relax
    \let\next\relax
  \else
    \ifnum\wordcount<3
      \ifnum\wordcount>0
        \edef\result{\result\space#1}%
      \else
        \edef\result{\result#1}%
      \fi
    \else
      \StrLeft{#1}{1}[\initial]%
      \edef\result{\result\space\initial.}%
    \fi
    \advance\wordcount by 1
    \def\next{\parsewords#2\relax}%
  \fi
  \next
}

\begin{document}
% Remove page number from first page
\thispagestyle{empty}

% Cover page
\begin{center}
    \makebox[\textwidth]{\includegraphics[width=0.4\textwidth]{img/Logo Isi.png}}\par
    \vspace{1cm}
    {\large\bfseries PROPOSAL PROGRAM KREATIVITAS MAHASISWA}\par
    \vspace{1cm}
    {\large\bfseries \judul}\par
    \vspace{1cm}

    {\large\bfseries BIDANG KEGIATAN\\PKM GAGASAN FUTURISTIK TERTULIS}
    \vspace{2cm}

    Disusun oleh:
    \vspace{0.5cm}

    \begin{tabular}{lcl}
        \expandafter\formatName\expandafter{\ketuaNama} & \ketuaNIM & Angkatan 2023\\
        \expandafter\formatName\expandafter{\anggotaSatuNama} & \anggotaSatuNIM & Angkatan 2024\\
        \expandafter\formatName\expandafter{\anggotaDuaNama} & \anggotaDuaNIM & Angkatan 2023\\
        \expandafter\formatName\expandafter{\anggotaTigaNama} & \anggotaTigaNIM & Angkatan 2024
    \end{tabular}
    \vspace{2cm}

    {\large\bfseries INSTITUT SENI INDONESIA\\YOGYAKARTA\\\the\year}
\end{center}

\newpage

% Start Roman numerals for front matter
\pagenumbering{roman}
\pagestyle{frontmatter}

% Table of contents
\renewcommand{\cftsecfont}{\normalfont}      % Normal font for section titles
\renewcommand{\cftsecpagefont}{\normalfont}  % Normal font for page numbers
\renewcommand{\cftsecleader}{\cftdotfill{\cftdotsep}} % Dots in TOC

\tableofcontents
\newpage

% Switch to Arabic numerals for main content
\pagenumbering{arabic}
\pagestyle{mainmatter}

% Main content
\section{PENDAHULUAN}

% Subsections can be added here based on specific requirements
\subsection{Latar Belakang}
% Add your background content here

\subsection{Rumusan Masalah}
% Add problem statements here

\subsection{Tujuan}
% Add objectives here

\subsection{Manfaat}
% Add benefits here

\newpage
\section{GAGASAN}
% Subsections can be added here based on specific requirements
\subsection{Kondisi Kekinian Pencetus Gagasan}
% \lipsum[1-3]
% Add current situation analysis here

\subsection{Solusi yang Pernah Ditawarkan}
% Add previous solutions here

\subsection{Gagasan Baru yang Ditawarkan}
% Add your new proposed solution here

\subsection{Pihak yang Dapat Mengimplementasikan Gagasan}
% Add stakeholders and implementation details here

\subsection{Langkah-langkah Strategis Implementasi Gagasan}
% Add strategic implementation steps here

\newpage
\section{KESIMPULAN}

% Add your conclusion content here
\subsection{Gagasan yang Diajukan}
% Summarize the main ideas proposed

\subsection{Teknik Implementasi}
% Summarize the implementation techniques

\subsection{Prediksi Hasil}
% Describe expected outcomes and impact

\newpage

% References and appendix
\section*{DAFTAR PUSTAKA}
% Use the following format for references:
% Nama penulis. Tahun. Judul. Tempat terbit: Penerbit.

% Example:
% Doe, John. 2023. Title of the Book. New York: Publisher Name.

\begin{flushleft}
% Add your references here in alphabetical order

\end{flushleft}

\newpage
\section{LAMPIRAN}

\subsection{Lampiran 1. Biodata Ketua dan Anggota, serta Dosen Pendamping}

{\setlength{\parindent}{0pt}
    \textbf{Biodata Ketua}\par
    \textbf{A. Identitas Diri}\par
    \vspace{6pt}
    \begin{tabularx}{\textwidth}{|c|l|
        >{\raggedright\arraybackslash\hspace{0pt}}X|}
    \hline
    1 & Nama Lengkap             & \ketuaNama   \\ \hline
    2 & Jenis Kelamin            & \ketuaGender \\ \hline
    3 & Program Studi            & \ketuaProdi  \\ \hline
    4 & NIM                      & \ketuaNIM    \\ \hline
    5 & Tempat dan Tanggal Lahir & \ketuaTTL    \\ \hline
    6 & Alamat E-mail            & \ketuaEmail  \\ \hline
    7 & Nomor Telepon/HP         & \ketuaTelp   \\ \hline
    \end{tabularx}\par

    \vspace{0.5cm}
    \textbf{B. Kegiatan Kemahasiswaan yang Sedang/Pernah Diikuti}\par
    \vspace{6pt}
    \begin{tabularx}{\textwidth}{|c|X|X|X|X|}
    \hline
    \multicolumn{1}{|c|}{\textbf{No}} & 
    \multicolumn{1}{>{\centering\arraybackslash}X|}{\makecell{\textbf{Jenis Kegiatan}}} & 
    \multicolumn{1}{>{\centering\arraybackslash}X|}{\makecell{\textbf{Status dalam}\\\textbf{Kegiatan}}} & 
    \multicolumn{1}{>{\centering\arraybackslash}X|}{\makecell{\textbf{Waktu dan Tempat}}} \\ \hline
    1 & Sarjana  & & \\ \hline
    2 & & & \\ \hline
    \end{tabularx}\par

    \vspace{0.5cm}
    \textbf{C. Penghargaan yang Pernah Diterima}\par
    \vspace{6pt}
    \begin{tabularx}{\textwidth}{|c|X|X|X|}
    \hline
    \multicolumn{1}{|c|}{\textbf{No}} & 
    \multicolumn{1}{>{\centering\arraybackslash}X|}{\makecell{\textbf{Jenis}\\\textbf{Penghargaan}}} & 
    \multicolumn{1}{>{\centering\arraybackslash}X|}{\makecell{\textbf{Pihak Pemberi}\\\textbf{Penghargaan}}} & 
    \multicolumn{1}{>{\centering\arraybackslash}X|}{\makecell{\textbf{Tahun}}} \\ \hline
    1 & & & \\ \hline
    2 & & & \\ \hline
    \end{tabularx}\par

    % Statement text at the end of each biodata
    \vspace{0.5cm}
    \sloppy Semua data yang saya isikan dan tercantum dalam biodata ini adalah benar dan dapat dipertanggungjawabkan secara hukum. Apabila di kemudian hari ternyata dijumpai ketidaksesuaian dengan kenyataan, saya sanggup menerima sanksi.

    \vspace{0.5cm}
    \sloppy Demikian biodata ini saya buat dengan sebenarnya untuk memenuhi salah satu persyaratan dalam pengajuan \textbf{PKM-GFT}.

    \vspace{1cm}
    \begin{flushright}
    Yogyakarta, \today\\
    Ketua Tim\\
    \vspace{2cm}
    (\ketuaNama)
    \end{flushright}
}
% Ketua to Anggota 1 ------------------------------

\newpage
{\setlength{\parindent}{0pt}
    \textbf{Biodata Anggota 1}\par
    \textbf{A. Identitas Diri}\par
    \vspace{6pt}
    \begin{tabularx}{\textwidth}{|c|l|
        >{\raggedright\arraybackslash\hspace{0pt}}X|}
    \hline
    1 & Nama Lengkap             & \anggotaSatuNama   \\ \hline
    2 & Jenis Kelamin            & \anggotaSatuGender \\ \hline
    3 & Program Studi            & \anggotaSatuProdi  \\ \hline
    4 & NIM                      & \anggotaSatuNIM    \\ \hline
    5 & Tempat dan Tanggal Lahir & \anggotaSatuTTL    \\ \hline
    6 & Alamat E-mail            & \anggotaSatuEmail  \\ \hline
    7 & Nomor Telepon/HP         & \anggotaSatuTelp   \\ \hline
    \end{tabularx}\par

    \vspace{0.5cm}
    \textbf{B. Kegiatan Kemahasiswaan yang Sedang/Pernah Diikuti}\par
    \vspace{6pt}
    \begin{tabularx}{\textwidth}{|c|X|X|X|}
    \hline
    \multicolumn{1}{|c|}{\textbf{No}} & 
    \multicolumn{1}{>{\centering\arraybackslash}X|}{\makecell{\textbf{Jenis Kegiatan}}} & 
    \multicolumn{1}{>{\centering\arraybackslash}X|}{\makecell{\textbf{Status dalam}\\\textbf{Kegiatan}}} & 
    \multicolumn{1}{>{\centering\arraybackslash}X|}{\makecell{\textbf{Waktu dan Tempat}}} \\ \hline
    1 & & & \\ \hline
    2 & & & \\ \hline
    \end{tabularx}\par

    \vspace{0.5cm}
    \textbf{C. Penghargaan yang Pernah Diterima}\par
    \vspace{6pt}
    \begin{tabularx}{\textwidth}{|c|X|X|X|}
    \hline
    \multicolumn{1}{|c|}{\textbf{No}} & 
    \multicolumn{1}{>{\centering\arraybackslash}X|}{\makecell{\textbf{Jenis}\\\textbf{Penghargaan}}} & 
    \multicolumn{1}{>{\centering\arraybackslash}X|}{\makecell{\textbf{Pihak Pemberi}\\\textbf{Penghargaan}}} & 
    \multicolumn{1}{>{\centering\arraybackslash}X|}{\makecell{\textbf{Tahun}}} \\ \hline
    1 & & & \\ \hline
    2 & & & \\ \hline
    \end{tabularx}\par

    % Statement text at the end of each biodata
    \vspace{0.5cm}
    \sloppy Semua data yang saya isikan dan tercantum dalam biodata ini adalah benar dan dapat dipertanggungjawabkan secara hukum. Apabila di kemudian hari ternyata dijumpai ketidaksesuaian dengan kenyataan, saya sanggup menerima sanksi.

    \vspace{0.5cm}
    \sloppy Demikian biodata ini saya buat dengan sebenarnya untuk memenuhi salah satu persyaratan dalam pengajuan \textbf{PKM-GFT}.

    \vspace{1cm}
    \begin{flushright}
    Yogyakarta, \today\\
    Anggota Tim\\
    \vspace{2cm}
    (\anggotaSatuNama)
    \end{flushright}
}

% Anggota 1 to Anggota 2 ------------------------------

\newpage
{\setlength{\parindent}{0pt}
    \textbf{Biodata Anggota 2}\par
    \textbf{A. Identitas Diri}\par
    \vspace{6pt}
    \begin{tabularx}{\textwidth}{|c|l|
        >{\raggedright\arraybackslash\hspace{0pt}}X|}
    \hline
    1 & Nama Lengkap             & \anggotaDuaNama   \\ \hline
    2 & Jenis Kelamin            & \anggotaDuaGender \\ \hline
    3 & Program Studi            & \anggotaDuaProdi  \\ \hline
    4 & NIM                      & \anggotaDuaNIM    \\ \hline
    5 & Tempat dan Tanggal Lahir & \anggotaDuaTTL    \\ \hline
    6 & Alamat E-mail            & \anggotaDuaEmail  \\ \hline
    7 & Nomor Telepon/HP         & \anggotaDuaTelp   \\ \hline
    \end{tabularx}\par

    \vspace{0.5cm}
    \textbf{B. Kegiatan Kemahasiswaan yang Sedang/Pernah Diikuti}\par
    \vspace{6pt}
    \begin{tabularx}{\textwidth}{|c|X|X|X|}
    \hline
    \multicolumn{1}{|c|}{\textbf{No}} & 
    \multicolumn{1}{>{\centering\arraybackslash}X|}{\makecell{\textbf{Jenis Kegiatan}}} & 
    \multicolumn{1}{>{\centering\arraybackslash}X|}{\makecell{\textbf{Status dalam}\\\textbf{Kegiatan}}} & 
    \multicolumn{1}{>{\centering\arraybackslash}X|}{\makecell{\textbf{Waktu dan Tempat}}} \\ \hline
    1 & & & \\ \hline
    2 & & & \\ \hline
    \end{tabularx}\par

    \vspace{0.5cm}
    \textbf{C. Penghargaan yang Pernah Diterima}\par
    \vspace{6pt}
    \begin{tabularx}{\textwidth}{|c|X|X|X|}
    \hline
    \multicolumn{1}{|c|}{\textbf{No}} & 
    \multicolumn{1}{>{\centering\arraybackslash}X|}{\makecell{\textbf{Jenis}\\\textbf{Penghargaan}}} & 
    \multicolumn{1}{>{\centering\arraybackslash}X|}{\makecell{\textbf{Pihak Pemberi}\\\textbf{Penghargaan}}} & 
    \multicolumn{1}{>{\centering\arraybackslash}X|}{\makecell{\textbf{Tahun}}} \\ \hline
    1 & & & \\ \hline
    2 & & & \\ \hline
    \end{tabularx}\par

    % Statement text at the end of each biodata
    \vspace{0.5cm}
    \sloppy Semua data yang saya isikan dan tercantum dalam biodata ini adalah benar dan dapat dipertanggungjawabkan secara hukum. Apabila di kemudian hari ternyata dijumpai ketidaksesuaian dengan kenyataan, saya sanggup menerima sanksi.

    \vspace{0.5cm}
    \sloppy Demikian biodata ini saya buat dengan sebenarnya untuk memenuhi salah satu persyaratan dalam pengajuan \textbf{PKM-GFT}.

    \vspace{1cm}
    \begin{flushright}
    Yogyakarta, \today\\
    Anggota Tim\\
    \vspace{2cm}
    (\anggotaDuaNama)
    \end{flushright}
}

% Anggota 2 to Anggota 3 ------------------------------

\newpage
{\setlength{\parindent}{0pt}
    \textbf{Biodata Anggota 2}\par
    \textbf{A. Identitas Diri}\par
    \vspace{6pt}
    \begin{tabularx}{\textwidth}{|c|l|
        >{\raggedright\arraybackslash\hspace{0pt}}X|}
    \hline
    1 & Nama Lengkap             & \anggotaTigaNama   \\ \hline
    2 & Jenis Kelamin            & \anggotaTigaGender \\ \hline
    3 & Program Studi            & \anggotaTigaProdi  \\ \hline
    4 & NIM                      & \anggotaTigaNIM    \\ \hline
    5 & Tempat dan Tanggal Lahir & \anggotaTigaTTL    \\ \hline
    6 & Alamat E-mail            & \anggotaTigaEmail  \\ \hline
    7 & Nomor Telepon/HP         & \anggotaTigaTelp   \\ \hline
    \end{tabularx}\par

    \vspace{0.5cm}
    \textbf{B. Kegiatan Kemahasiswaan yang Sedang/Pernah Diikuti}\par
    \vspace{6pt}
    \begin{tabularx}{\textwidth}{|c|X|X|X|}
    \hline
    \multicolumn{1}{|c|}{\textbf{No}} & 
    \multicolumn{1}{>{\centering\arraybackslash}X|}{\makecell{\textbf{Jenis Kegiatan}}} & 
    \multicolumn{1}{>{\centering\arraybackslash}X|}{\makecell{\textbf{Status dalam}\\\textbf{Kegiatan}}} & 
    \multicolumn{1}{>{\centering\arraybackslash}X|}{\makecell{\textbf{Waktu dan Tempat}}} \\ \hline
    1 & & & \\ \hline
    2 & & & \\ \hline
    \end{tabularx}\par

    \vspace{0.5cm}
    \textbf{C. Penghargaan yang Pernah Diterima}\par
    \vspace{6pt}
    \begin{tabularx}{\textwidth}{|c|X|X|X|}
    \hline
    \multicolumn{1}{|c|}{\textbf{No}} & 
    \multicolumn{1}{>{\centering\arraybackslash}X|}{\makecell{\textbf{Jenis}\\\textbf{Penghargaan}}} & 
    \multicolumn{1}{>{\centering\arraybackslash}X|}{\makecell{\textbf{Pihak Pemberi}\\\textbf{Penghargaan}}} & 
    \multicolumn{1}{>{\centering\arraybackslash}X|}{\makecell{\textbf{Tahun}}} \\ \hline
    1 & & & \\ \hline
    2 & & & \\ \hline
    \end{tabularx}\par

    % Statement text at the end of each biodata
    \vspace{0.5cm}
    \sloppy Semua data yang saya isikan dan tercantum dalam biodata ini adalah benar dan dapat dipertanggungjawabkan secara hukum. Apabila di kemudian hari ternyata dijumpai ketidaksesuaian dengan kenyataan, saya sanggup menerima sanksi.

    \vspace{0.5cm}
    \sloppy Demikian biodata ini saya buat dengan sebenarnya untuk memenuhi salah satu persyaratan dalam pengajuan \textbf{PKM-GFT}.

    \vspace{1cm}
    \begin{flushright}
    Yogyakarta, \today\\
    Anggota Tim\\
    \vspace{2cm}
    (\anggotaTigaNama)
    \end{flushright}
}

% Anggota 3 to Dosen ------------------------------

\newpage
{\setlength{\parindent}{0pt}
    \textbf{Biodata Dosen Pendamping}\par
    \textbf{A. Identitas Diri}\par
    \vspace{6pt}
    \begin{tabularx}{\textwidth}{|c|l|
        >{\raggedright\arraybackslash\hspace{0pt}}X|}
    \hline
    1 & Nama Lengkap (dengan gelar) & \dosenNama   \\ \hline
    2 & Jenis Kelamin               & \dosenGender \\ \hline
    3 & Program Studi               & \dosenProdi  \\ \hline
    4 & NIM                         & \dosenNIDN    \\ \hline
    5 & Tempat dan Tanggal Lahir    & \dosenTTL    \\ \hline
    6 & Alamat E-mail               & \dosenEmail  \\ \hline
    7 & Nomor Telepon/HP            & \dosenTelp   \\ \hline
    \end{tabularx}\par

    \vspace{0.5cm}
    \textbf{B. Riwayat Pendidikan}\par
    \vspace{6pt}
    \begin{tabularx}{\textwidth}{|c|X|X|c|c|}
    \hline
    \multicolumn{1}{|c|}{\textbf{No}} & 
    \multicolumn{1}{>{\centering\arraybackslash}X|}{\makecell{\textbf{Jenjang}}} & 
    \multicolumn{1}{>{\centering\arraybackslash}X|}{\makecell{\textbf{Bidang Ilmu}}} & 
    \multicolumn{1}{>{\centering\arraybackslash}X|}{\makecell{\textbf{Institusi}}} & 
    \multicolumn{1}{>{\centering\arraybackslash}X|}{\makecell{\textbf{Tahun Lulus}}} \\ \hline
    1 & Sarjana (S1) & Musik & ISI Yogyakarta & 1985\\ \hline
    2 & Magister (S2) & Humaniora/ Pengkajian Seni Pertunjukan & UGM & 1994 \\ \hline
    3 & Doktor (S3) & Humanioara/ Pengkajian Seni Pertunjukan dan Seni Rupa & UGM & 2010\\ \hline
    \end{tabularx}\par

    \vspace{0.5cm}
    \textbf{C. Rekam Jejak Tri Dharma PT (dalam 5 tahun terakhir)}\par
    \vspace{6pt}
    Pendidikan/Pengajaran\par
    \begin{tabularx}{\textwidth}{|c|X|l|l|}
        \hline
        \multicolumn{1}{|c|}{\textbf{No}} & 
        \multicolumn{1}{>{\centering\arraybackslash}X|}{\makecell{\textbf{Nama Mata Kuliah}}} & 
        \multicolumn{1}{>{\centering\arraybackslash}l|}{\makecell{\textbf{Wajib/Pilihan}}} & 
        \multicolumn{1}{>{\centering\arraybackslash}l|}{\makecell{\textbf{SKS}}} \\ \hline
        1 & Analisis Bentuk Musik & Wajib & 3\\ \hline
        2 & Kajian Literatur Musik & Wajib & 3\\ \hline
        3 & Musikologi III & Wajib & 2\\ \hline
    \end{tabularx}\par
    \vspace{0.5cm}
    
    Penelitian\par
    \begin{tabularx}{\textwidth}{|c|X|X|l|}
        \hline
        \multicolumn{1}{|c|}{\textbf{No}} & 
        \multicolumn{1}{>{\centering\arraybackslash}X|}{\makecell{\textbf{Judul Penelitian}}} & 
        \multicolumn{1}{>{\centering\arraybackslash}l|}{\makecell{\textbf{Penyandang Dana}}} & 
        \multicolumn{1}{>{\centering\arraybackslash}l|}{\makecell{\textbf{Tahun}}} \\ \hline
        1 & Kajian Historis Pembelajaran Gitar Klasik di Jurusan Musik FSP ISI Yogyakarta & ISI Yogyakarta& 2021\\ \hline
        2 & Kajian Musikologis Lantunan Surat al Fatihah oleh Imam Tetap Masjid Jogokariyan Yogyakarta & ISI Yogyakarta & 2022\\ \hline
        3 & Inovasi Asesmen Pembelajaran Kelas untuk Peningkatan Akurasi Penilaian Akhir Kuliah Analisis Bentuk Musik di Perguruan Tinggi Seni & ISI Yogyakarta & 2024\\ \hline
        4 & Hubungan antara Identitas Sosial dan Atribusi dengan Preferensi Musik pada Kalangan Remaja di Yogyakarta (anggota peneliti) & Kemdikbudristek/DRPM & 20204\\ \hline
    \end{tabularx}\par
    \vspace{0.5cm}

    Pengabdian Kepada Masyarakat\par
    \begin{tabularx}{\textwidth}{|c|X|X|l|}
        \hline
        \multicolumn{1}{|c|}{\textbf{No}} & 
        \multicolumn{1}{>{\centering\arraybackslash}X|}{\makecell{\textbf{Judul Pengabdian}\\\textbf{Kepada Masyarakat}}} & 
        \multicolumn{1}{>{\centering\arraybackslash}l|}{\makecell{\textbf{Penyandang Dana}}} & 
        \multicolumn{1}{>{\centering\arraybackslash}l|}{\makecell{\textbf{Tahun}}} \\ \hline
        1 & Editor Jurnal Promusika & ISI Yogyakarta& 2024\\ \hline
    \end{tabularx}\par
    \vspace{0.5cm}

    % Statement text at the end of each biodata
    \vspace{0.5cm}
    \sloppy Semua data yang saya isikan dan tercantum dalam biodata ini adalah benar dan dapat dipertanggungjawabkan secara hukum. Apabila di kemudian hari ternyata dijumpai ketidaksesuaian dengan kenyataan, saya sanggup menerima sanksi.

    \vspace{0.5cm}
    \sloppy Demikian biodata ini saya buat dengan sebenarnya untuk memenuhi salah satu persyaratan dalam pengajuan \textbf{PKM-GFT}.

    \vspace{1cm}
    \begin{flushright}
    Yogyakarta, \today\\
    Dosen Pendamping\\
    \vspace{2cm}
    (\dosenNama)
    \end{flushright}
}

\newpage
\subsection{Lampiran 2. Susunan Tim Pengusul dan Pembagian Tugas}
\begin{flushleft}
    \begin{tabularx}{\textwidth}{|c|X|X|X|X|}
    \hline
    \multicolumn{1}{|c|}{\textbf{No}} & 
    \multicolumn{1}{>{\centering\arraybackslash}X|}{\makecell{\textbf{Nama}}} & 
    \multicolumn{1}{>{\centering\arraybackslash}X|}{\makecell{\textbf{Posisi Penulis}}} & 
    \multicolumn{1}{>{\centering\arraybackslash}X|}{\makecell{\textbf{Bidang Ilmu}}} & 
    \multicolumn{1}{>{\centering\arraybackslash}X|}{\makecell{\textbf{Kontribusi}}} \\ \hline
    1 & \expandafter\formatName\expandafter{\ketuaNama}            &  & Musik & \\ \hline
    2 & \expandafter\formatName\expandafter{\anggotaSatuNama}      &  & Musik & \\ \hline
    3 & \expandafter\formatName\expandafter{\anggotaDuaNama}       &  & Musik & \\ \hline
    4 & \expandafter\formatName\expandafter{\anggotaTigaNama}      &  & Musik & \\ \hline
    5 & \dosenSingkatNama                                          &  & Musik & \\ \hline
    \end{tabularx}\par
\end{flushleft}

\newpage
\subsection{Lampiran 3. Surat Pernyataan Ketua Tim Pengusul}
{\setlength{\parindent}{0pt}
    \begin{center}
        \textbf{SURAT PERNYATAAN KETUA TIM PENGUSUL}
    \end{center}

    \vspace{0.5cm}
    Yang bertandatangan di bawah ini:

    \begin{tabularx}{\textwidth}{|l|c|X|}
        \hline
        Nama Ketua Tim & : & \ketuaNama \\ \hline
        Nomor Induk Mahasiswa & : & \ketuaNIM \\ \hline
        Program Studi & : & \ketuaProdi \\ \hline
        Nama Dosen pendamping & : & \dosenNama \\ \hline
        Perguruan Tinggi & : & Institut Seni Indonesia Yogyakarta \\ \hline
    \end{tabularx}

    \vspace{0.5cm}
    Dengan ini menyatakan bahwa proposal PKM-GFT saya dengan judul \judul{} yang diusulkan untuk tahun anggaran \the\year{} adalah:

    \begin{enumerate}[leftmargin=*]
        \item Asli karya kami, belum pernah dibiayai oleh lembaga atau sumber dana lain; dan\par
        \item Tidak dibuat dengan menggunakan kecerdasan buatan/\textit{artificial intelligence} (AI).\par
    \end{enumerate}

    \sloppy Bilamana di kemudian hari ditemukan ketidaksesuaian dengan pernyataan ini, maka saya bersedia dituntut dan diproses sesuai dengan ketentuan yang berlaku dan mengembalikan seluruh biaya yang sudah diterima ke kas negara.\par

    \vspace{0.3cm}
    \sloppy Demikian pernyataan ini dibuat dengan sesungguhnya dan dengan sebenar-benarnya.\par

    \vspace{1cm}
    \begin{tabular}{p{0.5\textwidth}p{0.5\textwidth}}
        & Yogyakarta, \today \\
        & Yang menyatakan, \\ [3cm]
        & (\ketuaNama) \\
        & NIM. \ketuaNIM \\
    \end{tabular}
}


\end{document}
